Since circulating an initial draft of this paper in March 2018,
datasheets for datasets have already gained traction in a number of
settings. Academic researchers have adopted our proposal and released
datasets with accompanying datasheets~\cite[e.g.,][]{recipeqa2018,
dataset1, choi2018quac, daume20gicoref}. Microsoft, Google, and IBM
have begun to pilot datasheets for datasets internally within product
teams. Researchers at Google published follow-up work on \emph{model
cards} that document machine learning models~\cite{model_cards} and
released a \emph{data card} (a lightweight version of a datasheet)
\edit{along} with the Open Images dataset~\cite{openimages}. Researchers at
IBM proposed \emph{factsheets}~\cite{factsheets} that document various
characteristics of AI services, including whether the datasets used to
develop the services are accompanied with datasheets. \edit{The Data
Nutrition Project incorporated some of the questions provided in the
previous section into the latest release of their Dataset Nutrition
Label}~\cite{nutritionsecondgen}. Finally, the Partnership on AI, a
multi-stakeholder organization focused on \edit{studying and
formulating} best practices for developing and deploying \edit{AI
technologies}, is working on industry-wide documentation guidance that
builds \edit{on} datasheets\edit{ for datasets}, model cards, and
factsheets.\footnote{\url{https://www.partnershiponai.org/about-ml/}}

These initial successes have also revealed implementation challenges
that may need to be addressed to support wider adoption. Chief among
them is the need for dataset creators to modify the questions and
workflow \edit{provided in the previous} section based on their existing
organizational infrastructure and workflows. We \edit{also note
that} \edit{the} questions and workflow
may \edit{pose} \edit{problems} for dynamic datasets. If a dataset
changes only infrequently, we recommend accompanying updated versions
with updated datasheets.

Datasheets for datasets do not provide a complete solution to
mitigating unwanted \edit{societal} biases or potential risks or
harms. Dataset creators cannot anticipate every possible use of a
dataset, and identifying unwanted \edit{societal} biases often requires
additional labels indicating demographic information \edit{about}
individuals, which may not be available to dataset creators for
reasons including those individuals' data protection and
privacy~\cite{holstein2018improving}.

When creating datasets that relate to people, \edit{and hence their
accompanying datasheets,} it may be necessary for dataset creators to
work with experts in other domains such as anthropology\edit{, sociology,
and science and technology studies.} There are complex and contextual
social, historical, and geographical factors that influence how best
to collect data \edit{from} individuals \edit{in a manner that is respectful}.

Finally, creating datasheets for datasets will necessarily impose
overhead on dataset creators. Although datasheets may reduce the
amount of time that dataset creators spend answering one-off questions
about datasets, the process of creating a datasheet will always take
time, and organizational infrastructure and workflows\edit{---not to mention
incentives---}will need to be modified to accommodate this investment.

\edit{Despite these implementation} challenges, there are many benefits to
creating datasheets for datasets. In addition to facilitating better
communication between dataset creators and dataset consumers,
datasheets provide \edit{an} opportunity for dataset creators to distinguish
themselves as prioritizing transparency and
accountability. Ultimately, we believe that the benefits to the
machine learning community outweigh the costs.
